\subsection{取向 torsor 的超对称结构、wreath 跳跃律与寄存器下界}\label{subsec:conclusion-orientation-torsor-supermonoidal-wreath-gauge-register}
沿定理 \ref{thm:xi-binary-jump-orientation-unification} 的统一口径，\(\mathrm{Or}\) 同时编码离散 uplift 与解析分支跳跃。
本小节将该口径在张量、乘积、分层与寄存器复杂度方向上的终态结构律统一为可直接调用的命题簇。

\begin{proposition}[\texorpdfstring{\(\ZZ_2\)}{Z2}-torsor 的平方范式化与奇偶幂坍缩]\label{prop:conclusion-z2torsor-square-parity-collapse}
设 \(P\) 为 \(\ZZ_2\)-torsor，并以
\[
P\otimes Q:=(P\times Q)/\Delta(\ZZ_2)
\]
定义 Picard 张量结构。
则存在不依赖额外选择的自然同构
\[
\Theta_P:\ P\otimes P\ \xrightarrow{\ \sim\ }\ \ZZ_2.
\]
从而对任意整数 \(n\ge 1\) 有规范同构
\[
P^{\otimes n}\cong
\begin{cases}
\ZZ_2,& n\ \text{为偶数},\\
P,& n\ \text{为奇数}.
\end{cases}
\]
\end{proposition}
\begin{proof}
任取 \((p,q)\in P\times P\)。
由 torsor 作用自由且传递，存在唯一 \(\varepsilon\in\ZZ_2\) 使 \(q=\varepsilon\cdot p\)。
定义
\[
\delta(p,q):=\varepsilon.
\]
对任意 \(\eta\in\ZZ_2\)，有 \(\delta(\eta\cdot p,\eta\cdot q)=\delta(p,q)\)，故 \(\delta\) 下推为
\[
\bar\delta:(P\times P)/\Delta(\ZZ_2)\to\ZZ_2.
\]
\(\bar\delta\) 显然双射，且与 \(\ZZ_2\)-作用相容，得 \(\Theta_P\)。
奇偶幂结论由 \(\Theta_P\) 迭代直接推出。证毕。
\end{proof}

\begin{theorem}[取向函子在不交并下的超对称单子结构]\label{thm:conclusion-orientation-supermonoidal-koszul}
对有限集 \(S,T\)，存在自然同构
\[
\mu_{S,T}:\ \mathrm{Or}(S)\otimes \mathrm{Or}(T)\xrightarrow{\ \sim\ }\mathrm{Or}(S\sqcup T),
\]
满足结合与幺元约束。
并且对交换同构 \(\tau_{S,T}:S\sqcup T\to T\sqcup S\) 有
\[
\tau_{S,T}^{\ast}\circ \mu_{S,T}
=
(-1)^{|S||T|}\,
\mu_{T,S}\circ \mathrm{swap}.
\]
等价地，\(\mathrm{Or}\) 在 \(\ZZ_2\)-分次意义下为超对称 monoidal，其 braiding 精确由 Koszul 因子 \((-1)^{|S||T|}\) 给出。
\end{theorem}
\begin{proof}
由命题 \ref{prop:bdry-orientation-functor-symmetric-monoidal-koszul} 可得对应结论（该命题采用 \(\mathrm{Or}(S\sqcup T)\to \mathrm{Or}(S)\otimes\mathrm{Or}(T)\) 的等价取向）。
将其取逆即得本式方向。
块交换奇偶性为 \((-1)^{|S||T|}\)，故得到 Koszul 因子。证毕。
\end{proof}

\begin{theorem}[笛卡尔积下的 Kronecker--取向分解]\label{thm:conclusion-orientation-cartesian-kronecker}
设有限集 \(S,T\) 的基数分别为 \(d,e\)。
则存在自然同构
\[
\Phi_{S,T}:\ \mathrm{Or}(S\times T)\xrightarrow{\ \sim\ }\mathrm{Or}(S)^{\otimes e}\otimes \mathrm{Or}(T)^{\otimes d}.
\]
\end{theorem}
\begin{proof}
命题 \ref{prop:bdry-orientation-cartesian-product-exponent-law} 已给出同一公式。
其证明本质上由行列式恒等式
\(
\det(U\otimes W)=\det(U)^{\dim W}\det(W)^{\dim U}
\)
与取向行列式实现（命题 \ref{prop:bdry-orientation-torsor-determinant-realization-jumpclass}）组成。
这正对应于标准 wreath 嵌入下的 Kronecker 分解机制。证毕。
\end{proof}

\begin{corollary}[积取向的二次坍缩与偶度屏蔽]\label{cor:conclusion-orientation-cartesian-parity-collapse}
在定理 \ref{thm:conclusion-orientation-cartesian-kronecker} 的同构下，结合命题 \ref{prop:conclusion-z2torsor-square-parity-collapse} 得
\[
\mathrm{Or}(S\times T)\ \cong\
\mathrm{Or}(S)^{\otimes (e\bmod 2)}
\otimes
\mathrm{Or}(T)^{\otimes (d\bmod 2)}.
\]
因此：
\begin{enumerate}
  \item \(d,e\) 皆偶时，\(\mathrm{Or}(S\times T)\cong \ZZ_2\)；
  \item \(d\) 偶、\(e\) 奇时，\(\mathrm{Or}(S\times T)\cong \mathrm{Or}(S)\)；
  \item \(d\) 奇、\(e\) 偶时，\(\mathrm{Or}(S\times T)\cong \mathrm{Or}(T)\)；
  \item \(d,e\) 皆奇时，\(\mathrm{Or}(S\times T)\cong \mathrm{Or}(S)\otimes\mathrm{Or}(T)\)。
\end{enumerate}
\end{corollary}
\begin{proof}
将定理 \ref{thm:conclusion-orientation-cartesian-kronecker} 中幂次按奇偶化简即可。证毕。
\end{proof}

\begin{theorem}[wreath 结构下的符号加法与偶度屏蔽]\label{thm:conclusion-wreath-sign-additivity-even-shielding}
令
\[
W:=\mathfrak S_d^{\,e}\rtimes \mathfrak S_e
\hookrightarrow
\mathfrak S_{de}
\]
为标准 wreath 嵌入，并把符号同态写作加法型
\(
\mathrm{sgn}_n:\mathfrak S_n\to\ZZ_2
\)。
则对任意 \(\big((\sigma_1,\dots,\sigma_e),\tau\big)\in W\) 有
\[
\mathrm{sgn}_{de}\!\big((\sigma_1,\dots,\sigma_e),\tau\big)
=
\Big(\sum_{j=1}^{e}\mathrm{sgn}_d(\sigma_j)\Big)+d\cdot \mathrm{sgn}_e(\tau)
\quad\text{于 }\ZZ_2.
\]
特别地：
\begin{enumerate}
  \item \(d\) 偶时，外层置换 \(\tau\) 对总二值跳跃不可见；
  \item \(d\) 奇时，总跳跃等于“块内跳跃之和 + 外层跳跃”。
\end{enumerate}
\end{theorem}
\begin{proof}
对应置换矩阵可分解为
\[
P=(P_\tau\otimes I_d)\cdot \mathrm{diag}(P_{\sigma_1},\dots,P_{\sigma_e}),
\]
故
\[
\det(P)=\det(P_\tau)^d\cdot\prod_{j=1}^{e}\det(P_{\sigma_j}).
\]
将行列式取 \(\ZZ_2\) 符号即得公式。两条推论是对 \(d\) 奇偶的直接代入。证毕。
\end{proof}

\begin{proposition}[fold-gauge 的中心二扭与二值电荷格分离]\label{prop:conclusion-foldbin-gauge-center-charge-gap}
设
\[
G_{\mathrm{bin},m}:=\mathcal{G}^{\mathrm{bin}}_m
=\prod_{x\in X_m}\mathfrak S_{d_m(x)},
\]
并定义
\[
s_m:=\#\{x\in X_m:\ d_m(x)\ge 2\},\qquad
t_m:=\#\{x\in X_m:\ d_m(x)=2\}.
\]
则有：
\begin{enumerate}
  \item \(\mathrm{Hom}(G_{\mathrm{bin},m},\ZZ_2)\cong (\ZZ_2)^{s_m}\)；
  \item \(Z(G_{\mathrm{bin},m})[2]\cong (\ZZ_2)^{t_m}\)；
  \item 存在自然坐标嵌入
  \[
  Z(G_{\mathrm{bin},m})[2]\hookrightarrow \mathrm{Hom}(G_{\mathrm{bin},m},\ZZ_2),
  \]
  其像恰为由二点纤维坐标生成的子空间；因此“非中心二值方向”的维数为 \(s_m-t_m\)。
\end{enumerate}
\end{proposition}
\begin{proof}
前两条分别由定理 \ref{thm:window6-foldbin-gauge-abelianization-even-parity} 与命题 \ref{prop:window6-foldbin-gauge-center-vs-charge-separation} 的结构分解直接给出。
嵌入由“二点纤维中心换位 \(\mapsto\) 该纤维符号字符”逐坐标定义；其单射性与像描述同样见命题 \ref{prop:window6-foldbin-gauge-center-vs-charge-separation} 的坐标化证明。
维数差即 \(s_m-t_m\)。证毕。
\end{proof}

\begin{theorem}[等变 \texorpdfstring{\(a\)}{a}-分裂的最小寄存器态下界]\label{thm:conclusion-equivariant-a-splitting-register-lower-bound}
令 \(\Delta\) 为 \(k\)-元集合，\(1\le a\le k-1\)，群 \(\mathfrak S_k\) 以自然方式作用于 \(\Delta\)。
称 \((R,A)\) 为一个 \(\mathfrak S_k\)-等变 \(a\)-分裂机制，若：
\begin{enumerate}
  \item \(R\) 是有限 \(\mathfrak S_k\)-集；
  \item \(A\subseteq \Delta\times R\) 在对角作用下不变；
  \item 对每个 \(r\in R\)，纤维 \(A_r:=\{x\in\Delta:(x,r)\in A\}\) 满足 \(|A_r|=a\)。
\end{enumerate}
则：
\begin{enumerate}
  \item 上述数据等价于一个 \(\mathfrak S_k\)-等变映射
  \[
  \varphi:R\to \binom{\Delta}{a},
  \qquad
  A=\{(x,r):x\in\varphi(r)\};
  \]
  \item 必有
  \[
  |R|\ge \binom{k}{a};
  \]
  \item 下界可达：取 \(R=\binom{\Delta}{a}\) 与 \(\varphi=\mathrm{id}\) 即得。
\end{enumerate}
特别地，当 \(a=1\) 时，任一“\(1\sqcup(k-1)\)”等变选层机制都需至少 \(k\) 个寄存器态，即至少 \(\log_2 k\) 比特。
\end{theorem}
\begin{proof}
由 \(|A_r|=a\) 可定义 \(\varphi(r):=A_r\in\binom{\Delta}{a}\)。
\(A\) 的不变性等价于
\(
\varphi(\sigma\cdot r)=\sigma\cdot\varphi(r)
\)，故 \(\varphi\) 等变，且 \(A\) 由 \(\varphi\) 唯一恢复，得第 (1) 条。

\(\binom{\Delta}{a}\) 在 \(\mathfrak S_k\) 作用下传递，故非空等变像必为全体 \(\binom{\Delta}{a}\)。
于是
\(
|R|\ge |\mathrm{Im}\,\varphi|=\binom{k}{a}
\)，得第 (2) 条。
第 (3) 条由恒等映射构造直接成立。
最后一句取 \(a=1\) 即得。证毕。
\end{proof}

\begin{theorem}[奇基数非平凡二值跳跃对超对称的结构强迫]\label{thm:conclusion-odd-cardinality-nontrivial-jump-forces-super-symmetry}
设
\[
F:\mathbf{FinSet}^{\cong}\to \mathbf{Tors}(\ZZ_2)
\]
为 monoidal 函子，源端张量为不交并 \(\sqcup\)，目标端张量为 torsor 张量 \(\otimes\)。
若要求其与标准交换同构严格按“无额外符号”的对称约束相容，则对任意奇整数 \(d\ge 3\)，限制
\[
F|_{\mathbf{Set}^{\cong}_d}
\]
必为平凡函子。
等价地：一旦在某个奇基数层出现非平凡二值跳跃（由定理 \ref{thm:bdry-z2-jump-functor-uniqueness} 必同构于 \(\mathrm{Or}\)），交换约束必引入 Koszul 因子，故体系只能是超对称 monoidal，而非无符号对称 monoidal。
\end{theorem}
\begin{proof}
反设存在奇 \(d\ge 3\) 使 \(F|_{\mathbf{Set}^{\cong}_d}\) 非平凡。
由定理 \ref{thm:bdry-z2-jump-functor-uniqueness}，该限制对应符号类。
取 \(|S|=|T|=d\)。
则在块对角子群 \(\mathfrak S_d\times\mathfrak S_d\subset \mathfrak S_{2d}\) 上，\(F\) 的字符限制为两块符号之和；故 \(F\) 在 \(\mathfrak S_{2d}\) 上亦为非平凡字符。
由唯一性只能是 \(\mathrm{sgn}_{2d}\)。
对块交换 \(\beta_{S,T}\)，其奇偶为 \((-1)^{d^2}=-1\)，因此 \(F(\beta_{S,T})\) 必非平凡。

另一方面，若交换约束无额外符号，则当 \(S=T\) 时，\(F(\beta_{S,S})\) 在
\(
F(S)\otimes F(S)
\)
上仅为对换。
由命题 \ref{prop:conclusion-z2torsor-square-parity-collapse} 的规范同构
\(
F(S)\otimes F(S)\cong \ZZ_2
\)，该对换作用为恒等，故 \(F(\beta_{S,S})\) 平凡，矛盾。
故假设不成立。证毕。
\end{proof}

